\subsection{Macaulay 模与有限代数}

注.--- 设 \(\rho: A \to B\) 为一个环同态，\(\mathfrak{p} \in \operatorname{Spec}(A)\)。用 \(\overline{B}\) 表示环 \(\kappa(\mathfrak{p}) \otimes_{A} B\)。它等同于 \(S^{-1}B/\mathfrak{p}(S^{-1}B)\)，其中 \(S\) 是形如 \(\rho(A - \mathfrak{p})\) 的 \(B\) 的乘性子集；因此 \(\overline{B}\) 的素理想就是诸理想 \(\mathfrak{q}\overline{B}\)，其中 \(\mathfrak{q}\) 是 \(B\) 的包含 \(\mathfrak{p}B\) 且不与 \(S\) 相交的素理想，换言之，即 \(B\) 的位于 \(\mathfrak{p}\) 之上的素理想。对这样的理想 \(\mathfrak{q}\)，有 \(S \subset B - \mathfrak{q}\)，故 \(\overline{B}\) 在 \(\mathfrak{q}\overline{B}\) 处的局部环等同于 \(B_{\mathfrak{q}}/pB_{\mathfrak{q}}\)，亦即 \(\kappa(\mathfrak{p}) \otimes_{A} B_{\mathfrak{q}}\)。

类似地，若 N 是一个 B-模，则 \(\overline{\mathrm{B}}_{\mathfrak{q}\overline{\mathrm{B}}}\) -模 \(\left(\kappa(\mathfrak{p})\otimes_{\mathrm{A}}\mathrm{N}\right)_{\mathfrak{q}\overline{\mathrm{B}}}\) 等同于 \(\kappa(\mathfrak{p})\otimes_{\mathrm{A}}\mathrm{N}_{\mathfrak{q}}\)。再设 B-模 N 是有限生成的；由 Nakayama 引理，条件 \(\kappa(\mathfrak{p})\otimes_{\mathrm{A}}\mathrm{N}_{\mathfrak{q}}=0\) 等价于 \(\mathrm{N}_{\mathfrak{q}}=0\)。于是 \(\overline{\mathrm{B}}\) -模 \(\kappa(\mathfrak{p})\otimes_{\mathrm{A}}\mathrm{N}\) 的支撑集由诸理想 \(\mathfrak{q}\overline{\mathrm{B}}\) 组成，其中 \(\mathfrak{q}\) 取遍 \(\operatorname{Supp}_{\mathrm{B}}(\mathrm{N})\) 中位于 \(\mathfrak{p}\) 之上的素理想。特别地，为使模 \(\kappa(\mathfrak{p})\otimes_{\mathrm{A}}\mathrm{N}\) 非零，必须且只需存在 \(\operatorname{Supp}_{\mathrm{B}}(\mathrm{N})\) 的一个位于 \(\mathfrak{p}\) 之上的素理想。

命题 8.--- 设 \(\rho: A \to B\) 为 Noether 环之间的同态，\(N\) 为一个 \(B\) -模，它作为 \(A\) -模是有限型的。为使 \(A\) -模 \(N\) 是 Macaulay 的，必须且只需 \(B\) -模 \(N\) 是 Macaulay 的，并且对每对 \((\mathfrak{n}, \mathfrak{n}')\)（\(\mathrm{Supp}_{\mathrm{B}}(N)\) 中满足 \(\rho^{-1}(\mathfrak{n}) = \rho^{-1}(\mathfrak{n}')\) 的极大理想），有 \(\dim_{B_{\mathfrak{n}}} (N_{\mathfrak{n}}) = \dim_{B_{\mathfrak{n}}'} (N_{\mathfrak{n}}')\)。

A-模 B/Ann \(_B(N)\) 同构于有限生成 A-模 \(\text{End}_A(N)\) 的一个子模，因此是有限型的。把 A 换成 A/Ann \(_A(N)\)，B 换成 B/Ann \(_B(N)\)，我们就归结到下述情形：\(\rho\) 是单射且使 B 成为有限 A-代数，并且有 \(\text{Supp}_A(N) = \text{Spec}(A)\) 及 \(\text{Supp}_B(N) = \text{Spec}(B)\)。映射 \(f: \text{Spec}(B) \to \text{Spec}(A)\)（由 \(\rho\) 诱导）这时是满射，且 B 的素理想 \(\mathfrak{q}\) 是极大的当且仅当 \(f(\mathfrak{q})\) 是 A 的极大理想 (V, § 2, 第 1 节, 定理 1 及命题 1)。

设 m 为 A 的一个极大理想。由上述注，环 \(B_{m}\) 的包含 \(mB_{m}\) 的素理想是形如 \(qB_{m}\) 的理想，其中 \(q \in \text{Spec}(B)\) 是 B 的（必然为极大的）理想，且满足 \(f(\mathfrak{q}) = \mathfrak{m}\)。我们有

\[
\operatorname{prof} _ {\mathrm{A} _ {\mathfrak {m}}} \left(\mathrm{N} _ {\mathfrak {m}}\right) = \operatorname{prof} _ {\mathrm{B} _ {\mathfrak {m}}} \left(\mathfrak {m B} _ {\mathfrak {m}}; \mathrm{N} _ {\mathfrak {m}}\right) = \inf _ {\mathfrak {q} \in f ^ {- 1} (\mathfrak {m})} \left(\operatorname{prof} _ {\mathrm{B} _ {\mathfrak {q}}} \left(\mathrm{N} _ {\mathfrak {q}}\right)\right)
\]

(§ 1, 第 3 节, 命题 4 及第 5 节, 命题 8)，以及

\[
\dim_ {A _ {\mathfrak {m}}} \left(N _ {\mathfrak {m}}\right) = \dim_ {B _ {\mathfrak {m}}} \left(N _ {\mathfrak {m}}\right) = \sup _ {\mathfrak {q} \in f ^ {- 1} (\mathfrak {m})} \left(\dim_ {B _ {\mathfrak {q}}} \left(N _ {\mathfrak {q}}\right)\right)
\]

(VIII, § 2, 第 3 节, 定理 1 及 § 1, 第 4 节, 命题 9)。由于有 \(\mathrm{prof}_{\mathcal{B}_{\mathfrak{q}}}(\mathbf{N}_{\mathfrak{q}}) \leqslant \dim_{\mathcal{B}_{\mathfrak{q}}}(\mathbf{N}_{\mathfrak{q}})\)（对每个 \(\mathfrak{q} \in f^{-1}(\mathfrak{m})\)），本命题由这些等式得出。

推论 1.--- 设 \(\rho: A \to B\) 为 Noether 环之间的同态。若 B 是有限 A-代数且是 Macaulay A-模，则它是 Macaulay 环。此外，若 \(\rho\) 是单射，则对 A 的每个理想 a 有 ht(aB) = ht(a)，且对 B 的每个理想 b 有 ht(b) = ht(ρ \(^{-1}\) (b))。

第一个断言由命题 8 得出。设 \(\rho\) 是单射。设 \(\mathfrak{a}\) 为 A 的理想。我们有 \(\mathrm{ht}(\mathfrak{a}) = \mathrm{prof}_{\mathrm{A}}(\mathfrak{a};\mathrm{B})\)，因为 A-模 B 是 Macaulay 的且其支撑集等于 \(\mathrm{Spec}(\mathrm{A})\) (\(n^{\circ}1\), 命题 1 的推论)；又 \(\mathrm{ht}(\mathfrak{a}\mathrm{B}) = \mathrm{prof}_{\mathrm{B}}(\mathfrak{a}\mathrm{B};\mathrm{B})\)（同上引文），且 \(\mathrm{prof}_{\mathrm{A}}(\mathfrak{a};\mathrm{B}) = \mathrm{prof}_{\mathrm{B}}(\mathfrak{a}\mathrm{B};\mathrm{B})\) (\(\S 1, n^{\circ}3\), 命题 4)，由此得 \(\mathrm{ht}(\mathfrak{a}\mathrm{B}) = \mathrm{ht}(\mathfrak{a})\)。设 \(\mathfrak{b}\) 为 B 的理想。由前所述，有 \(\mathrm{ht}(\boldsymbol{\rho}^{-1}(\mathfrak{b})) = \mathrm{ht}(\boldsymbol{\rho}^{-1}(\mathfrak{b})\mathrm{B})\)。但 \(\boldsymbol{\rho}^{-1}(\mathfrak{b})\mathrm{B}\) 含于 \(\mathfrak{b}\)，故其高度小于 \(\mathrm{ht}(\mathfrak{b})\)，而有 \(\mathrm{ht}(\mathfrak{b}) \leqslant \mathrm{ht}(\boldsymbol{\rho}^{-1}(\mathfrak{b}))\)（由 VIII, \(\S 2, n^{\circ}3\), 定理 1, b)）。

推论 2.--- 设 A 为整闭的 Noether 环，B 为包含 A 的环。假设 B 作为 A-模是有限型的无扭模。若 B 是 Macaulay 环，则 A-模 B 是 Macaulay 的。

事实上，B 中位于 A 的同一理想之上的两个素理想具有相同的高度 (VIII, § 2, 第 3 节, 定理 2)。因此可以对 N = B 应用命题 8。

推论 3.--- 设 A 为整闭环，K 为其分式域，L 为一个有限 K-代数，使得 {[}L : K{]} \(1_{\mathrm{A}}\) 在 A 中可逆，B 为 L 的一个子 A-代数，且在 A 上有限。

\begin{enumerate}
\def\labelenumi{\alph{enumi})}
\item
  B 的子 A-模 A1B 是一个直和因子。
\item
  对 A 的每个理想 J，有不等式 \(\operatorname{prof}_{\mathrm{A}}(\mathrm{J};\mathrm{A})\geqslant\operatorname{prof}_{\mathrm{B}}(\mathrm{JB};\mathrm{B})\)。
\item
  若 B 是 Macaulay 环，则 A 亦然。
\end{enumerate}

K-线性映射 \(\mathrm{Tr}_{\mathrm{L/K}}: \mathrm{L} \to \mathrm{K}\) 把 B 映到 A 中 (V, § 1, 第 6 节, 命题 17 的推论 2)，因而通过限制定义了一个 A-线性映射 \(t: \mathrm{B} \to \mathrm{A}\)。对每个 \(x \in \mathrm{A}\)，有 \(t(x1_{\mathrm{B}}) = [\mathrm{L}: \mathrm{K}]x\)，由此得到 a)。

由 § 1, 第 3 节, 命题 4，有 \(\operatorname{prof}_{\mathrm{A}}(\mathrm{J};\mathrm{B}) = \operatorname{prof}_{\mathrm{B}}(\mathrm{JB};\mathrm{B})\)；而由 a) 及 § 1, 第 1 节, 注 4，有 \(\operatorname{prof}_{\mathrm{A}}(\mathrm{J};\mathrm{A}) \geqslant \operatorname{prof}_{\mathrm{A}}(\mathrm{J};\mathrm{B})\)，由此得到 b)。

若环 B 是 Noether 的，则 A 亦然：事实上，由 a)，有 \(\mathfrak{a}\mathrm{B} \cap \mathrm{A} = \mathfrak{a}\)（对 A 的每个理想 \(\mathfrak{a}\)）；于是 A 的每个递增理想序列 \((\mathfrak{a}_n)_{n \in \mathbf{N}}\) 都是稳定的，因为序列 \((\mathfrak{a}_n\mathrm{B})_{n \in \mathbf{N}}\) 是稳定的。在 c) 的假设下，A-模 B 是 Macaulay 的（推论 2），从而 A-模 A 亦然 (第 1 节, 例 2)。

例.--- 推论 3 特别适用于下述两种情形：

\begin{enumerate}
\def\labelenumi{\alph{enumi})}
\item
  取一个 Noether 整闭环 A、其分式域的一个有限次 n 可分扩张 L（n 满足 \(n1_{A}\) 在 A 中可逆），并取 B 为 A 在 L 中的整闭包 (V, § 1, 第 6 节, 命题 18 的推论 1)。
\item
  取一个 Noether 整闭环 B 以及 B 的自同构的一个有限群 G，使得 \(\mathrm{Card(G)1_B}\) 在 B 中可逆。我们取 A 为 B 中在 G 的作用下不变的元素所组成的环。我们来验证这是 a) 的一个特殊情形。群 G 作用在 B 的分式域 L 上，而 L 对这个作用的不变元域就是 A 的分式域 K (V, § 1, 第 9 节, 命题 23 的推论)。L 是 K 的 Galois 扩张，从而更是可分扩张；其 Galois 群同构于 G，故 {[}L : K{]} 等于 Card G。{[}L : K{]} \(_{1_{B}}\) 的逆在 G 下不变，故 {[}L : K{]} \(_{1_{A}}\) 在 A 中可逆。由于 B 是整闭的，环 A（即 K ∩ B）是整闭的，且 B 是 A 在 L 中的整闭包（同上引文, 命题 22）。
\end{enumerate}

特别地，若 B 是 Macaulay 环，则 A 亦然。

